% Cuadrícula de 5 mm (A4) — Eric Lucero González — https://ericlucero.dev
% Licencia CC BY-NC-SA 4.0. Compilar con pdfLaTeX (dos pasadas: latexmk lo hace solo).
\documentclass[a4paper]{article}
\usepackage[T1]{fontenc}
\usepackage[spanish]{babel}
\usepackage{lmodern}
\usepackage{tikz}
\pagestyle{empty}

% ── Configuración (mm) ─────────────────────────────────────
\def\anchoPapel{210}
\def\altoPapel{297}
\def\margenMin{5}   % margen mínimo (las impresoras no llegan al borde)
\def\paso{5}        % lado de cada cuadro

\begin{document}
\begin{tikzpicture}[
    remember picture, overlay, shift={(current page.south west)},
    x=1mm, y=1mm,
    grid lines/.style={gray!50, very thin}
  ]
  % Número de cuadros enteros que caben, y la cuadrícula centrada en la hoja
  \pgfmathtruncatemacro{\nx}{(\anchoPapel - 2*\margenMin) / \paso}
  \pgfmathtruncatemacro{\ny}{(\altoPapel  - 2*\margenMin) / \paso}
  \pgfmathsetmacro{\ancho}{\nx * \paso}
  \pgfmathsetmacro{\alto}{\ny * \paso}
  \pgfmathsetmacro{\x}{(\anchoPapel - \ancho) / 2}
  \pgfmathsetmacro{\y}{(\altoPapel  - \alto) / 2}

  \draw[step=\paso, grid lines, shift={(\x,\y)}] (0,0) grid (\ancho,\alto);
\end{tikzpicture}
\end{document}
